\section{Discussion}
\label{sec:discussion}

\subsection{Why Learning Rate Controlled Everything}

The results show a very clear pattern: the learning rate was the single most important factor in whether a model succeeded or failed. At $10^{-6}$, almost every model struggled. This makes sense when you consider that we kept the base convolutional layers frozen. Since we were starting the classification head from scratch, those weights needed a decent-sized push to start learning the features extracted by the pre-trained model. At $10^{-6}$, the updates were so tiny that the model barely moved from its initial random state before our \texttt{EarlyStopping} callback cut it off.

It is worth noting that early stopping any of the configurations does not confirm whether or not the respective model had finished learning. Note that the callback was set with a patience of 10 epochs and a minimum delta of 0.001 on \texttt{val\_auc}, indicating that the run was halted only when there was no meaningful improvement over 10 consecutive epochs. Within the 72 runs conducted in this study, this limit was set to build consistency amongst the results over seeking out peak performance. Thus, it is worth acknowledging that the results for 1e-6 configurations are representative of the learning rate's performance under the specific conditions of the study as opposed to the upper limit of what those configuraitons were capable of achieving given more time. 
\subsection{Comparing the Optimizers}

A notable observation was how the different optimizers handled the relatively small dataset. RMSProp was the winner here, followed closely by Adam and Adamax. These adaptive methods are great for medical imaging projects like this because they can adjust their updates based on the gradients they see, which helps when the data is limited and potentially noisy.

In contrast, AdaDelta was surprisingly poor. It is meant to be an improvement over Adagrad, but on the provided dataset, it seemed to stall almost immediately. This may be due to, with so few steps per epoch, AdaDelta's internal learning rate calculation dropping to zero before the model could find a good path. 

\subsection{Architecture Strengths}

ResNet50 was the most reliable performer across the arhitecture suite. This may be due to its skip connections helping keep the gradient signal clean as it passes back to the classification head. VGG19 actually rendered the single best individual run, which might be because its simpler, older architecture produces very straightforward features that the pooling-and-dense head could easily interpret. 

EfficientNetB3, which is the most modern model tested, actually came in last on average. Which may be surprising, but it's possible that its more complex feature space needs a lot more data than our 518 training images to really perform to specification.

\subsection{A Warning about AUC}

One of the most important takeaways for us was about the AUC metric itself. We saw several suboptimal runs where the model predicted the same class for every single image, yet they still had an AUC over 0.80. For instance, the VGG19 + Adam (LR=$10^{-6}$) (Figure~\ref{fig:cm-vgg19-adam-lr1e6}) with gradient updates so miniscule, the model weights never meaningully left from start; leading to all 247 test images to be predicted as cancer without exception, an F1 of 0, and a test accuracy of 48.58\% (seemingly random chance) despite a retention of AUC 0.5520. This happens because AUC measures how well the model ranks images, not how well it classifies them at a specific threshold. In a real clinical setting, a model that never finds a single cancer case is a total failure, even if its ranking ability looks good on paper. This reinforces the need to always check recall and specificity alongside AUC.

\subsection{Recommendations}

Should the research be continued, the next logical step would be fine-tuning. We kept the base models frozen to see what a clean transfer learning performance looked like, but unfreezing the last few layers and training them at a very low learning rate would likely boost the scores significantly. It would be prudent to try $k$-fold cross-validation to get a better sense of how these models generalize across different subsets of the data. Furthermore, future efforts are recommended to conduct follow-up simulations with a higher early stopping threshold (e.g. 20 or 30 more epochs), or no early stopping at all. The current simulation setup would not be able to meaningfully distinguish between a configuration that failed to learn anything, and a configuration that was terminated before it was given any meaningful time to learn. Isolating learning rate would further clarify whether the learning rate is reflective of the true upper limits of configuration performance or whether it is an artifact of the stopping criteria. 
